\documentclass{article}
\usepackage[utf8]{inputenc}
\usepackage[utf8]{inputenc}
\usepackage{graphicx}
\usepackage{subcaption}
\usepackage{float}
\usepackage{amsmath}
\usepackage{amssymb}
\usepackage{enumerate}
\usepackage{physics}
\usepackage{siunitx}
\usepackage{tikz}
\usepackage{circuitikz}
\usepackage{indentfirst}

\addtolength{\oddsidemargin}{-.875in}
\addtolength{\evensidemargin}{-.875in}
\addtolength{\textwidth}{1.75in}

\addtolength{\topmargin}{-.875in}
\addtolength{\textheight}{1.75in} 
\title{HWtemplate}
\author{shomikchatterjee }
\date{April 2021}

\begin{document}
\begin{center}
\textbf{
{\Large ECE 210 Midterm 3 Worksheet}
}
\end{center} 
\noindent\makebox[\linewidth]{\rule{\linewidth}{0.4pt}}

\paragraph{Note:} This worksheet is not guaranteed to be entirely representative of the midterm's contents. Material may appear on this worksheet which will not appear on the midterm, and vice versa.

\subsection*{Partial Fraction Expansion}

\paragraph{1)} (Technique 1)

\subparagraph{a)} sub problem 1

\subparagraph{b)} sub problem 2


\vfill

\paragraph{2)} (Technique 2)

\subparagraph{a)} sub problem 1

\subparagraph{b)} sub problem 2


\vfill

\newpage
\subsection*{S-Domain Circuit Analysis}

\paragraph{3)} Consider the circuit below (SAMPLE CIRCUIT - CHANGE).

\begin{figure}[ht!]
\centering
\begin{circuitikz}[american, transform shape, voltage dir = old]
\draw (-3,0) to [R=$3\Omega$] (-3,3) to[short] (0,3) to[short] (3,3);
\draw (3,0) to [cI,l=$\alpha v_x$] (3,3);
\draw (0,3) to[R=4$\Omega$, v_=$v_x$] (0,0);
\draw (-3,0) to[short] (3,0);
\draw (-3,-3) to[V,l=$2\text{V}$] (-3,0) to[short] (3,0);
\draw (0,-3) to [R = 3$\Omega$] (0,0);
\draw (-3,-3) to[R=1$\Omega$] (-1.5,-3) to[V,l=1V] (0,-3);
\draw (0,-3) to[R=2$\Omega$] (3,-3) to[short](3,0);
\draw[-latex] (-3,-3.5) -- (-1.5,-3.5);
\node at (-2.25, -3.5) [anchor=north] {$i_t$};
\end{circuitikz}
\end{figure}

\subparagraph{a)} sub problem 1

\subparagraph{b)} sub problem 2

\vfill

\newpage

\subsection*{Marginal Stability \& Resonance}

\paragraph{4)} problem 1

\subparagraph{a)} marginal stability part

\subparagraph{b)} resonance part

\subparagraph{c)} etc

\vfill

\paragraph{5)} problem 2

\subparagraph{a)} marginal stability part

\subparagraph{b)} resonance part

\subparagraph{c)} etc

\vfill


\end{document}
